\section{The Appointed Texts}
\label{sec:appointed-texts}

The Missal's texts for this Sunday are those of \latin{Dominica I Adventus} in
the \work{Missale Romanum}, editio typica tertia (2002), which the
\work{Roman Missal}, Third Edition, renders for the dioceses of the United
States; the readings, psalm and Alleluia verse are those of no.~2 of the
\work{Lectionary for Mass}, which follows the \work{Ordo lectionum Missae} of
1981. The Latin of the Missal, the approved English of the Missal and the
approved English of the Lectionary are all under copyright, and none of them
is reproduced here. Each biblical text is given at its verses in the
Douay--Rheims version as revised by Bishop Challoner, which translates the
Clementine Vulgate. That English is a study text and is not the wording
proclaimed at Mass. Words printed in brackets and italics belong to the verse
or passage but not to the appointed text. Each prayer is identified by the
opening words of its Latin and described, and no English is supplied for it.
The two antiphons adapt Scripture: the verses they draw on are printed beside
them as their sources and not as their words, and no English is composed for
either.

\Needspace{8\baselineskip}
\subsection{Entrance Antiphon}
\label{proper-entrance-antiphon}

\properlocus{\latin{Ad te levavi animam meam}. The Missal cites Psalm 24:1--3
in the Vulgate numbering, which is Psalm 25:1--3 in the Hebrew, and prints
\latin{Cf.}}

\noindent The speaker lifts up his soul to God and says that he trusts in him;
he asks not to be put to shame, nor mocked by his enemies, for none who wait
for God will be confounded. The antiphon follows the first three verses of the
psalm with three departures. It has no vocative after \latin{Ad te}, where
both the Clementine Vulgate and the \work{Nova Vulgata} have \latin{Domine}.
It keeps the Clementine's \latin{Neque irrideant me inimici mei}, ``Neither
let my enemies laugh at me'', where the \work{Nova Vulgata} has another verb.
And for the waiting of the last line it reads \latin{qui te exspectant},
``who wait for thee'', a phrase found in neither Latin psalter: the Clementine
and the \work{Nova Vulgata} have \latin{qui sustinent te}. No English is
printed for the antiphon. The verses it draws on follow.

\begin{englishwitness}{Douay--Rheims (Challoner), Ps 24:1--3, the antiphon's source}
\textsuperscript{1}\notread{Unto the end, a psalm for David.} To thee, \notread{O
Lord,} have I lifted up my soul. \textsuperscript{2}In thee, O my God, I put
my trust; let me not be ashamed. \textsuperscript{3}Neither let my enemies
laugh at me: for none of them that wait on thee shall be confounded.
\end{englishwitness}

\Needspace{8\baselineskip}
\subsection{Collect}
\label{proper-collect}

\properlocus{\latin{Da, quaesumus, omnipotens Deus, hanc tuis fidelibus
voluntatem}. Long conclusion.}

\noindent Addressed to the almighty God, a petition for a will: that he give
his faithful the will (\latin{voluntatem}) to go out to meet his Christ as he
comes (\latin{Christo tuo venienti}) with just works (\latin{iustis operibus
occurrentes}), so that, set at his right hand, they may be found worthy to
possess the kingdom of heaven. No English is printed.

\Needspace{8\baselineskip}
\subsection{First Reading}
\label{proper-first-reading}

\properlocus{Isaiah 63:16b--17, 19b; 64:2b--7, in the Lectionary's numbering,
which follows the Hebrew; Lectionary no.~2. The Douay--Rheims, like the
Clementine Vulgate it translates, divides the chapters differently: the cry of
63:19b is its 64:1, and the Lectionary's 64:2--7 is its 64:3--8. The reading
is in three parts: the second half of 63:16 and all of 63:17; the cry of
63:19b; and 64:2b--7. The \work{Ordo}'s Latin begins \latin{Tu, Domine, pater
noster}.}

\begin{englishwitness}{Douay--Rheims (Challoner), Is 63:16--17; 64:1, 3--8}
\textsuperscript{16}\notread{For thou art our father, and Abraham hath not
known us, and Israel hath been ignorant of us:} thou, O Lord, art our father,
our redeemer, from everlasting is thy name. \textsuperscript{17}Why hast thou
made us to err, O Lord, from thy ways: why hast thou hardened our heart, that
we should not fear thee? return for the sake of thy servants, the tribes of
thy inheritance.

\notread{63:18--19 and 64:2 are not read.}

\textsuperscript{64:1}O that thou wouldst rend the heavens, and wouldst come
down: the mountains would melt away at thy presence.

\textsuperscript{3}When thou shalt do wonderful things, we shall not bear
them: thou didst come down, and at thy presence the mountains melted away.
\textsuperscript{4}From the beginning of the world they have not heard, nor
perceived with the ears: the eye hath not seen, O God, besides thee, what
things thou hast prepared for them that wait for thee. \textsuperscript{5}Thou
hast met him that rejoiceth, and doth justice: in thy ways they shall remember
thee: behold thou art angry, and we have sinned: in them we have been always,
and we shall be saved. \textsuperscript{6}And we are all become as one
unclean, and all our justices as the rag of a menstruous woman: and we have
all fallen as a leaf, and our iniquities, like the wind, have taken us away.
\textsuperscript{7}There is none that calleth upon thy name: that riseth up,
and taketh hold of thee: thou hast hid thy face from us, and hast crushed us
in the hand of our iniquity. \textsuperscript{8}And now, O Lord, thou art our
father, and we are clay: and thou art our maker, and we all are the works of
thy hands.
\end{englishwitness}

\noindent The Douay verse 64:3 is printed whole because the editions read
different halves of it. It has two members: ``When thou shalt do wonderful
things, we shall not bear them'' and ``thou didst come down, and at thy
presence the mountains melted away''. The \work{Ordo}'s Latin begins at the
second, \latin{Descendisti, et a facie tua montes defluxerunt}, so that in
Latin the mountains flow twice, once in the cry and once as a thing done. The
American calendar and page of readings cite the whole verse as 64:2, and the
page's English has
the clause about the wonderful deeds that were not hoped for, the first
member, and nothing for the second; the bishops' current translation of the
Bible notes that it omits words the Hebrew repeats there from 63:19. Where a
claim below rests on \latin{Descendisti}, it rests on the Latin.

\Needspace{8\baselineskip}
\subsection{Responsorial Psalm}
\label{proper-responsorial-psalm}

\properlocus{Psalm 80 (79):2a and 2c, 3b, 15--16, 18--19; the response is
v.~4. Lectionary no.~2. The second member of v.~2, the first member of v.~3,
vv.~5--14 and v.~17 are not sung, and the title is not part of the first
strophe. The \work{Ordo}'s response is \latin{Deus, converte nos, illustra
faciem tuam, et salvi erimus}.}

\begin{englishwitness}{Douay--Rheims (Challoner), Ps 79:1--4, 15--19}
\textsuperscript{1}\notread{Unto the end, for them that shall be changed, a
testimony for Asaph, a psalm.} \textsuperscript{2}Give ear, O thou that rulest
Israel: \notread{thou that leadest Joseph like a sheep.} Thou that sittest
upon the cherubims, shine forth \textsuperscript{3}\notread{Before Ephraim,
Benjamin, and Manasses.} Stir up thy might, and come to save us.
\textsuperscript{4}Convert us, O God: and shew us thy face, and we shall be
saved.

\notread{vv.~5--14 are not sung.}

\textsuperscript{15}Turn again, O God of hosts, look down from heaven, and
see, and visit this vineyard: \textsuperscript{16}And perfect the same which
thy right hand hath planted: and upon the son of man whom thou hast confirmed
for thyself. \notread{\textsuperscript{17}Things set on fire and dug down shall
perish at the rebuke of thy countenance.}
\textsuperscript{18}Let thy hand be upon the man of thy right hand: and upon
the son of man whom thou hast confirmed for thyself.
\textsuperscript{19}And we depart not from thee, thou shalt quicken us: and we
will call upon thy name.
\end{englishwitness}

\noindent Verse~4 is the refrain of the whole psalm; it returns, a little
varied, as vv.~8 and~20. The \work{Ordo}'s response asks God to make his face
shine (\latin{illustra}), the word of the \work{Nova Vulgata}; the Clementine
and the Douay--Rheims have ``shew'' (\latin{ostende}). At v.~16 the Latin
psalters and the Douay--Rheims end the verse with the son of man, as v.~18
ends; a version translated from the Hebrew, the King James, has there ``the
branch that thou madest strong for thyself''. The second strophe proclaimed
in the United States ends with the son of man, as the Latin does.

\Needspace{8\baselineskip}
\subsection{Second Reading}
\label{proper-second-reading}

\properlocus{1~Corinthians 1:3--9. Lectionary no.~2. The \work{Ordo}'s Latin
begins \latin{Fratres: Gratia vobis et pax}.}

\begin{englishwitness}{Douay--Rheims (Challoner), 1~Cor 1:3--9}
\textsuperscript{3}Grace to you and peace, from God our father and from the
Lord Jesus Christ. \textsuperscript{4}I give thanks to my God always for you,
for the grace of God that is given you in Christ Jesus:
\textsuperscript{5}That in all things you are made rich in him, in all
utterance and in all knowledge; \textsuperscript{6}As the testimony of Christ
was confirmed in you, \textsuperscript{7}So that nothing is wanting to you in
any grace, waiting for the manifestation of our Lord Jesus Christ.
\textsuperscript{8}Who also will confirm you unto the end without crime, in
the day of the coming of our Lord Jesus Christ. \textsuperscript{9}God is
faithful: by whom you are called unto the fellowship of his Son, Jesus Christ
our Lord.
\end{englishwitness}

\noindent ``Without crime'' renders the Clementine's \latin{sine crimine}; the
Greek word means one who cannot be accused, and the Lectionary's English
renders the sense, not the Latin word. ``Manifestation'' renders
\latin{revelationem}, the revelation of Christ.

\Needspace{8\baselineskip}
\subsection{Alleluia Verse}
\label{proper-acclamation}

\properlocus{Psalm 85 (84):8. Lectionary no.~2, printed without \latin{Cf.}
The \work{Ordo} gives the verse of Year~A again: \latin{Ostende nobis, Domine,
misericordiam tuam, et salutare tuum da nobis}, word for word the Clementine
Vulgate.}

\begin{englishwitness}{Douay--Rheims (Challoner), Ps 84:8}
\textsuperscript{8}Shew us, O Lord, thy mercy; and grant us thy salvation.
\end{englishwitness}

\Needspace{8\baselineskip}
\subsection{Gospel}
\label{proper-gospel}

\properlocus{Mark 13:33--37. Lectionary no.~2. The \work{Ordo} opens the
reading \latin{In illo tempore: Dixit Iesus discipulis suis: Videte,
vigilate}. Verse~32, printed here in brackets, stands just before the reading
and is not read.}

\begin{englishwitness}{Douay--Rheims (Challoner), Mk 13:32--37}
\notread{\textsuperscript{32}But of that day or hour no man knoweth, neither
the angels in heaven, nor the Son, but the Father.}
\textsuperscript{33}Take ye heed, watch and pray. For ye know not when the
time is. \textsuperscript{34}Even as a man who, going into a far country, left
his house and gave authority to his servants over every work and commanded the
porter to watch. \textsuperscript{35}Watch ye therefore (for you know not
when the lord of the house cometh, at even, or at midnight, or at the cock
crowing, or in the morning): \textsuperscript{36}Lest coming on a sudden, he
find you sleeping. \textsuperscript{37}And what I say to you, I say to all:
Watch.
\end{englishwitness}

\noindent The Douay--Rheims ``and pray'' in v.~33 translates the Clementine
Vulgate's \latin{et orate}; the \work{Ordo}'s opening of the reading,
\latin{Videte, vigilate}, does not show those words. The porter is the
doorkeeper (\latin{ianitor}) of the house. The four times of v.~35 are the
four watches into which the night was divided.

\Needspace{8\baselineskip}
\subsection{Prayer over the Offerings}
\label{proper-prayer-over-offerings}

\properlocus{\latin{Suscipe, quaesumus, Domine, munera quae de tuis offerimus
collata beneficiis}. Short conclusion.}

\noindent A petition in two members: that the Lord accept the gifts his people
offer from what he has bestowed on them (\latin{de tuis \ldots\ collata
beneficiis}), and that what he grants their devotion to accomplish in this
passing time (\latin{temporali}) may become for them the reward of eternal
redemption. No English is printed.

\Needspace{8\baselineskip}
\subsection{Preface}
\label{proper-preface}

\properlocus{Preface I of Advent, \latin{De duobus adventibus Christi}, ``On
the two comings of Christ''. Appointed by the Missal's rubric for this Mass;
the Preface's own rubric assigns it to the days of Advent from the First
Sunday to 16~December. The thanksgiving proper to it begins \latin{Qui, primo
adventu in humilitate carnis assumptae}.}

\noindent Thanks are given to the Father through Christ, who at his first
coming, in the lowliness of the flesh he had taken, fulfilled the design
prepared long before and opened the way of salvation, so that when he comes a
second time in the glory of his majesty, those who now keep watch and dare to
wait (\latin{vigilantes nunc audemus exspectare}) may receive openly what was
promised. No English is printed.

\Needspace{8\baselineskip}
\subsection{Communion Antiphon}
\label{proper-communion-antiphon}

\properlocus{\latin{Dominus dabit benignitatem}. The Missal cites Psalm 84:13
in the Vulgate numbering, which is Psalm 85:13 in the Hebrew, and prints no
\latin{Cf.}}

\noindent A promise of the Lord's gift and of the land's harvest. The
antiphon is the whole verse without its first word, \latin{Etenim},
``For''. No English is printed for the antiphon; its source follows, with the
word it omits in brackets.

\begin{englishwitness}{Douay--Rheims (Challoner), Ps 84:13, the antiphon's source}
\textsuperscript{13}\notread{For} the Lord will give goodness: and our earth
shall yield her fruit.
\end{englishwitness}

\Needspace{8\baselineskip}
\subsection{Prayer after Communion}
\label{proper-prayer-after-communion}

\properlocus{\latin{Prosint nobis, quaesumus, Domine, frequentata mysteria}.
Short conclusion.}

\noindent A petition that the mysteries in which the assembly has shared may
profit it, with the reason: by them God teaches his people even now, while
they live among things that pass away (\latin{inter praetereuntia}), to love
the things of heaven and to cling to the things that will last
(\latin{inhaerere mansuris}). No English is printed.

\Needspace{8\baselineskip}
\subsection{Solemn Blessing}
\label{proper-solemn-blessing}

\properlocus{\latin{Omnipotens et misericors Deus, cuius Unigeniti adventum}.
The solemn blessing for Advent, which the Missal's rubric says may be used at
the end of this Mass; it is optional.}

\noindent Three invocations, each answered \latin{Amen}, followed by the
Trinitarian blessing. The first asks God to sanctify those who believe that
his Only-begotten has come and who await his coming again; the second, to
make them firm in faith, glad in hope and active in charity; the third, that
those who rejoice in the Redeemer's coming in the flesh may receive eternal
life when he comes again in majesty. No English is printed.
